% =========================================================
% Lab 1 — Basic Statistics (statement)
% Time Series Analysis — Chapter 1
% Author : Aymen Ben Brik (aymen.benbrik@esprit.tn)
% =========================================================
\documentclass[a4paper,11pt]{article}

\newcommand{\TPnumero}{Lab 1 -- Basic Statistics}
% =========================================================
% Tutorial / lab preamble — Time Series Analysis module
% article class + fancyhdr + tcolorbox + listings (English)
% Author : Aymen Ben Brik (aymen.benbrik@esprit.tn)
% =========================================================

\usepackage[utf8]{inputenc}
\usepackage[T1]{fontenc}
\usepackage[english]{babel}
\usepackage{lmodern}
\usepackage{amsmath, amssymb, amsthm, mathtools, bm}
\usepackage{graphicx}
\usepackage{booktabs}
\usepackage{array}
\usepackage{multirow}
\usepackage{colortbl}
\usepackage{xcolor}
\usepackage{tikz}
\usetikzlibrary{positioning, arrows.meta, calc, shapes, fit, backgrounds, matrix}
\usepackage{listings}
\usepackage[most]{tcolorbox}
\usepackage{geometry}
\geometry{a4paper, margin=2cm, headheight=15pt}
\usepackage{fancyhdr}
\usepackage[hidelinks]{hyperref}
\usepackage{enumitem}

% --- Colours ---
\definecolor{espritBlue}{RGB}{30, 60, 110}
\definecolor{espritAccent}{RGB}{200, 80, 30}
\definecolor{boxEx}{RGB}{30, 60, 110}
\definecolor{boxQ}{RGB}{180, 120, 0}
\definecolor{boxC}{RGB}{30, 100, 60}
\definecolor{boxAtt}{RGB}{200, 80, 30}

% --- Listings ---
\definecolor{codebg}{RGB}{248,248,248}
\definecolor{codeframe}{RGB}{220,220,220}
\definecolor{keyword}{RGB}{0,82,155}
\definecolor{comment}{RGB}{88,110,117}
\definecolor{string}{RGB}{133,153,0}

\lstdefinestyle{pythonstyle}{
  language=Python,
  basicstyle=\ttfamily\small,
  keywordstyle=\color{keyword}\bfseries,
  commentstyle=\color{comment}\itshape,
  stringstyle=\color{string},
  backgroundcolor=\color{codebg},
  frame=single,
  rulecolor=\color{codeframe},
  frameround=tttt,
  breaklines=true,
  showstringspaces=false,
  tabsize=2
}
\lstdefinestyle{Rstyle}{
  language=R,
  basicstyle=\ttfamily\small,
  keywordstyle=\color{keyword}\bfseries,
  commentstyle=\color{comment}\itshape,
  stringstyle=\color{string},
  backgroundcolor=\color{codebg},
  frame=single,
  rulecolor=\color{codeframe},
  frameround=tttt,
  breaklines=true,
  showstringspaces=false,
  tabsize=2
}
\lstset{style=pythonstyle}

% --- Common macros (configurable path) ---
\providecommand{\macrosPath}{../../preamble/macros.tex}
\input{\macrosPath}

% --- Exercise counter ---
\newcounter{excounter}
\renewcommand{\theexcounter}{\arabic{excounter}}

\newtcolorbox[use counter=excounter]{exercise}[1][]{
  enhanced, breakable,
  colback=boxEx!5, colframe=boxEx, fonttitle=\bfseries,
  title={Exercise~\theexcounter\ifstrempty{#1}{}{ -- #1}},
  arc=2pt, boxrule=0.6pt
}
\newtcolorbox{question}[1][]{
  enhanced, breakable,
  colback=boxQ!5, colframe=boxQ, fonttitle=\bfseries,
  title={Question\ifstrempty{#1}{}{ -- #1}},
  arc=2pt, boxrule=0.5pt
}
\newtcolorbox{solution}[1][]{
  enhanced, breakable,
  colback=boxC!5, colframe=boxC, fonttitle=\bfseries,
  title={Solution\ifstrempty{#1}{}{ -- #1}},
  arc=2pt, boxrule=0.5pt
}
\newtcolorbox{warning}[1][]{
  enhanced, breakable,
  colback=boxAtt!5, colframe=boxAtt, fonttitle=\bfseries,
  title={Warning\ifstrempty{#1}{}{ -- #1}},
  arc=2pt, boxrule=0.5pt
}

% --- Headers / footers ---
\pagestyle{fancy}
\fancyhf{}
\fancyhead[L]{\textsc{\moduleNom} -- \auteurNom}
\fancyhead[R]{\TPnumero}
\fancyfoot[C]{\thepage}
\fancyfoot[L]{\auteurInstitution}
\fancyfoot[R]{\auteurEmail}
\renewcommand{\headrulewidth}{0.4pt}
\renewcommand{\footrulewidth}{0.2pt}

% --- Placeholder ---
\providecommand{\TPnumero}{Lab}


\title{Lab 1: Basic statistics in Python\\
\large Moments $\cdot$ MoM $\cdot$ MLE $\cdot$ Inference}
\author{\auteurNom\\\auteurEmail\\\auteurInstitution}
\date{\anneeUniversitaire}

\begin{document}
\maketitle

\section*{Goals}
\begin{itemize}
  \item Compute empirical moments and compare with theoretical values.
  \item Implement Method of Moments and MLE for the Gaussian.
  \item Study the sampling distribution of the sample mean by
        Monte-Carlo simulation.
  \item Build a confidence interval and run a $t$-test on a real
        dataset (Atlanta monthly temperatures).
\end{itemize}

\section*{Setup}
\begin{itemize}
  \item Python 3, dependencies: \texttt{numpy}, \texttt{scipy},
        \texttt{matplotlib}.
  \item Reproducibility: \texttt{numpy.random.seed(42)}.
  \item Dataset: \texttt{data/AvTempAtlanta.txt} — monthly
        average temperatures, Atlanta, USA.
  \item The skeleton \texttt{seed.py} contains function signatures
        and the visualisation utilities.
\end{itemize}

% =========================================================
\begin{exercise}[Empirical moments]
\textbf{(1)} Generate $n = 10\,000$ samples from $\mathcal{N}(2, 4)$
(mean $2$, variance $4$). Compute the four empirical moments:
\[
  \widehat{\mu}, \quad \widehat{\sigma}^2,
  \quad \widehat{S} = \tfrac{1}{n\,\widehat\sigma^3}\sum (X_i - \widehat\mu)^3,
  \quad \widehat{K} = \tfrac{1}{n\,\widehat\sigma^4}\sum (X_i - \widehat\mu)^4.
\]
Compare with theoretical values $(\mu, \sigma^2, S, K) = (2, 4, 0, 3)$.

\textbf{(2)} Repeat for $n = 100$. Plot $\widehat{\mu}$ as a function of
$n$ for $n \in \{10, 50, 100, 500, 1000, 5000, 10\,000\}$.
What do you observe?
\end{exercise}

% =========================================================
\begin{exercise}[Method of Moments vs MLE for the Gaussian]
\textbf{(1)} Implement \texttt{mom\_gaussian(x)} and
\texttt{mle\_gaussian(x)} returning the estimated $(\widehat\mu,
\widehat\sigma^2)$. Verify numerically that the two estimators
\emph{coincide} on the same sample.

\textbf{(2)} Compare the MLE/MoM estimator $\widehat\sigma^2_{\text{MLE}}
= \tfrac{1}{n}\sum (X_i - \bar X)^2$ with the unbiased one
$S^2_n = \tfrac{1}{n-1}\sum (X_i - \bar X)^2$. Run $B = 1000$ Monte-Carlo
replicates of size $n = 20$, draw the histograms of both estimators,
and verify empirically that
\[
  \E\bigl[\widehat\sigma^2_{\text{MLE}}\bigr] \;=\; \tfrac{n-1}{n}\,\sigma^2,
  \qquad
  \E\bigl[S^2_n\bigr] \;=\; \sigma^2.
\]
\end{exercise}

% =========================================================
\begin{exercise}[Sampling distribution of $\bar X_n$]
\textbf{(1)} For $n \in \{10, 30, 100, 1000\}$, simulate
$B = 1000$ independent samples $X_1, \dots, X_n \iid
\mathcal{N}(\mu, \sigma^2) = \mathcal{N}(2, 4)$, and compute the
sample mean of each.

\textbf{(2)} Plot four histograms (one per $n$) of the $B$ values of
$\bar X_n$. Overlay the theoretical density of
$\mathcal{N}(\mu, \sigma^2/n)$. Compute the empirical variance and
compare with $\sigma^2 / n$.

\textbf{(3)} Repeat with a non-Gaussian distribution
($\mathrm{Exp}(\lambda = 1)$, mean and variance both equal to $1$).
What does the CLT predict, and is it observed?
\end{exercise}

% =========================================================
\begin{exercise}[Confidence interval and $t$-test on Atlanta temperatures]
Load \texttt{data/AvTempAtlanta.txt} (one column of monthly
temperatures, $n$ values).

\textbf{(1)} Compute $\bar X_n$ and $S_n$. Build the 95\% $t$-confidence
interval for $\mu$.

\textbf{(2)} Test $H_0 : \mu = 60$ versus $H_1 : \mu \neq 60$ at level
$5\%$. Report the test statistic, the critical value, the $p$-value,
and the conclusion.

\textbf{(3)} Compare with \texttt{scipy.stats.ttest\_1samp}.
\end{exercise}

% =========================================================
\begin{exercise}[Bonus: bootstrap CI]
\textbf{(1)} Implement a non-parametric bootstrap CI for $\mu$ (resample
the data $B = 5000$ times with replacement, compute $\bar X^*$ on each
resample, take the empirical $(\alpha/2, 1 - \alpha/2)$-quantiles of
the bootstrap distribution).

\textbf{(2)} Compare the bootstrap CI with the $t$-CI on the Atlanta
data. Discuss when each is preferable.
\end{exercise}

\bigskip
\begin{warning}
Always set the seed (\texttt{np.random.seed(42)}) before any random
draw. Save figures to PNG with explicit names
(\texttt{lab1\_ex3\_sampling.png}, etc.).
\end{warning}

\end{document}
